\chapter{Toolchain Checklist}
\label{ch:app-b-toolchain}

This appendix records how to rebuild the opening chapter examples and the book document without guessing at tools or commands, covering the language baseline, warning discipline, lab build commands, sanitizer practice, document toolchain, word counting, and partial and complete builds.
A reader who follows these steps should obtain the same programs and document structure the manuscript describes, so treat any mismatch as evidence rather than a detail to work around silently.

\section{Language baseline}
\label{sec:app-b-baseline}

All example C code uses the C11 language and the C standard library only, with no external C dependency.
Host intermediate values reach $510$ for addition and $65025$ for multiplication, which both fit before the explicit remainder operation because C11 guarantees that \code{unsigned int} represents at least the range $0$ through $65535$~\cite[Section~5.2.4.2.1]{wg14-2011-c11}.
The code therefore does not depend on \code{unsigned} having any particular width, on an optional exact-width type, or on a host overflow accidentally producing the guest language result.
Constants intended as unsigned carry the \code{U} suffix, as in \code{256U}.
\index{C11}

\section{Strict warnings}
\label{sec:app-b-warnings}

Every lab command builds with \code{-std=c11 -Wall -Wextra -Wpedantic -Werror}.
The standard flag pins the language version, the three warning flags report a broad set of suspicious constructs, and the final flag turns those warnings into errors so a build with a new warning cannot pass silently.
Both of the routinely used compilers must accept the code cleanly under this combination, so a warning from either compiler is a defect in the example rather than a difference of opinion to suppress locally.
Do not add a warning-suppression flag, lower the standard level, or remove \code{-Werror} to accommodate untested code without understanding the underlying construct.
If a new example genuinely needs a different option, record the reason in its README entry and keep the default strict set for everything else.

\section{Lab build commands}
\label{sec:app-b-lab-commands}

Run all commands from the \code{example-code/ch01/} directory and place executables outside the manuscript tree, under a temporary directory.
The expression, demo, and observation labs build as follows.
\begin{lstlisting}[caption={Build and run the expression, demo, and observation labs.},label={lst:app-b-basic-labs}]
cc -std=c11 -Wall -Wextra -Wpedantic -Werror expr.c test_expr.c -o /tmp/opencode/ch01-test
/tmp/opencode/ch01-test
cc -std=c11 -Wall -Wextra -Wpedantic -Werror expr.c demo.c -o /tmp/opencode/ch01-demo
/tmp/opencode/ch01-demo
cc -std=c11 -Wall -Wextra -Wpedantic -Werror observations.c -o /tmp/opencode/ch01-observations
/tmp/opencode/ch01-observations
\end{lstlisting}
The first program checks both binary operations on all operand pairs, the worked example on all inputs, repeated folding, generated trees, and the arithmetic exercise answers, while the second prints selected before-and-after values for inspection against the chapter tables.
The third program is a separate complete C program whose expected output is \code{value=0 calls=1}.
\begin{lstlisting}[caption={Build and run the stack-code lab.},label={lst:app-b-stack-lab}]
cc -std=c11 -Wall -Wextra -Wpedantic -Werror expr.c stack.c test_stack.c -o /tmp/opencode/ch01-stack
/tmp/opencode/ch01-stack
\end{lstlisting}
This program checks literal operand pairs for both stack arithmetic instructions, the exact worked-example code, peak stack demand before and after folding, generated-tree translation, malformed-code rejection, and capacity rollback.
\begin{lstlisting}[caption={Build and run the fault and trap laboratories.},label={lst:app-b-fault-trap-labs}]
cc -std=c11 -Wall -Wextra -Wpedantic -Werror expr.c stack.c fault_lab.c -o /tmp/opencode/ch01-fault
/tmp/opencode/ch01-fault
cc -std=c11 -Wall -Wextra -Wpedantic -Werror trap_lab.c -o /tmp/opencode/ch01-trap
/tmp/opencode/ch01-trap
\end{lstlisting}
The fault lab mutates emitted \code{MUL8} instructions into \code{ADD8} as an isolated specimen, while the trap lab uses its own expression type with division, an event log, and fuel-limited loops.
Keep assertions enabled when running every test program, since disabled assertions remove the executable form of the input and ownership contracts.

\section{Sanitizers and assertions}
\label{sec:app-b-sanitizers}

Passing a strict-warning build is necessary but not sufficient evidence, so the chapter verification also runs every test program under a \term{address sanitizer} and undefined-behavior sanitizers with both routinely used compilers.
The sanitizer runs cover the expression tests, the stack tests, and the fault and trap laboratories.
AddressSanitizer detects many invalid accesses, and UndefinedBehaviorSanitizer detects selected undefined operations on executed paths.
Neither is a general detector of uninitialized-value reads; those need separate instrumentation such as MemorySanitizer where supported.
A clean run means those checks reported no failure on the tested executions, not that the program is free of defects.
Assertions stay enabled during all of these runs and document preconditions such as non-null child pointers, valid node tags, initialized code objects, non-null output pointers, and inputs belonging to $V$.
Do not mutate a code object between verification and execution, and do not feed the executor uninitialized slots or an oversized logical length, since those actions step outside the verified interface.

\section{Document toolchain}
\label{sec:app-b-document-build}

The book is typeset with XeLaTeX by default or LuaLaTeX as an alternative.
The pdfLaTeX engine will not work because the document class uses \code{fontspec}.
The bibliography is processed with Biber through the \code{biblatex} backend and the index with MakeIndex.
The build driver is \code{build.lua} and it runs under \code{texlua}.
The required system fonts are Source Serif 4 for body text, Source Sans 3 for headings and captions, and JetBrains Mono for code.
Listing~\ref{lst:app-b-document-builds} summarizes the driver commands.
\begin{lstlisting}[caption={Document build, check, and cleanup commands.},label={lst:app-b-document-builds}]
texlua build.lua
texlua build.lua draft 3 4 7
texlua build.lua quick
texlua build.lua check
texlua build.lua fonts
texlua build.lua clean
texlua build.lua distclean
ENGINE=lualatex texlua build.lua
\end{lstlisting}
The full build runs the engine, Biber, MakeIndex, and two further engine passes, while a draft build compiles only named chapters for speed and a quick build performs a single engine pass without bibliography or index processing.
The check command verifies that expected source files exist.
The fonts command verifies that the required fonts are installed.
The clean command removes auxiliary files and the distclean command additionally removes the generated PDF.
The \code{ENGINE} environment setting selects LuaLaTeX instead of the default XeLaTeX.
Never commit generated auxiliary files, bibliographies, index files, synchronization data, or the PDF unless the maintainer explicitly asks.

\section{Word-count command}
\label{sec:app-b-word-count}

Chapter length is audited with TeXcount rather than estimated by eye.
The agreed command for the opening chapter is \code{texcount -inc -opt=support/texcount.opts chapters/ch01.tex}.
The \code{-inc} option includes files pulled in by the chapter, and the reported sum weights prose, headings, and captions.
Executable code, pseudocode, mathematical symbols, bibliography entries, and index entries are excluded from the prose count.
There is no chapter-length quota; coverage and clarity determine length.
Record the TeXcount command and its tool version together with every word audit, since a bare number without its command cannot be compared with later audits.
A broken exclusion can inflate or deflate every chapter count without changing any chapter source.

\section{Partial and complete builds}
\label{sec:app-b-partial-complete}

A partial build compiles only existing or named chapters, so cross-references into omitted chapters appear as \enquote{??}, which is expected draft behavior rather than a defect.
A complete build compiles the entire book with all chapters and appendices present, and only it can validate every cross-chapter reference, citation, and the full index.
All 38 chapter files and all three appendices are present, so the complete-book source check can pass.
Source presence and successful typesetting do not establish technical correctness or completion of the implementation exercises.
Promote a chapter toward complete only when its partial build is clean, meaning no undefined references, no undefined citations, no duplicate labels, and no package warnings.
Inspect rendered figures, tables, code listings, bibliography entries, and index entries in the PDF rather than trusting the log alone, and treat overfull boxes wider than a few points as defects to fix.
Run the full build before declaring document-wide work finished, including changes to the class file, the style file, the master document, citations, or the index.
